% page 440 of the facsimile (image p-439 of the scan)
% block 157.5 x 246.3 mm ; left margin 8.0 mm
\pgc{-12.700mm}{0.000mm}{
\l{\cel{2}432}
\l{}
\l{\sou{pepsino}. (\sou{kemio).}\cel{2}Substanco efikiva di la suko stomakala,}
\l{\cel{3}qua dissolvas la materii albuminoida. - DEFIRS.}
\l{}
\l{\sou{peptono}.\cel{2}(\sou{farmacio}).Produkturo de la digesto artificala}
\l{\cel{3}di karno da pepsino o da pankreatino, en kondicioni de-}
\l{\cel{3}terminita di medio e di temperaturo. - DEFIRS.}
\l{}
\l{\sou{per.}\cel{3}Prepoziciono qua indikas la instrumento di la ago,}
\l{\cel{3}to quo uzesas por produktar lu, la moyeno.}
\l{}
\l{\sou{pero.}\cel{2}Vasalo suprega-ranga, membro di la judicio-korto}
\l{\cel{3}rejala. - DEFIRS.}
\l{}
\l{\sou{perceptar}. - (\sou{trans.})\cel{2}Kaptar, komprenar, la donataji de}
\l{\cel{3}la sensi. - DEFIS.}
\l{}
\l{\sou{perchar}. (\sou{netrans.})\cel{2}(\sou{aludante la uceli}). Okupar kustume}
\l{\cel{3}plaso sur brancho di arboro, sur stango, sur paliso.-}
\l{\cel{3}EFS.}
\l{}
\l{\sou{perdar}. (\sou{trans)}.\cel{2}Des-aquirar profitajo, havajo. - II.}
\l{\cel{3}Cesar havar. - FIS.}
\l{}
\l{\sou{perdriko}.\cel{2}(\sou{zool.)}\cel{2}Ucelo galinacea, ne-domestika, qua vi-}
\l{\cel{3}vas kun sua simili en la frument-agri, en la prati.-FIS}
\l{}
\l{\sou{perena}.\cel{2}(\sou{bot.})\cel{2}Qua ri-kreskas dum plura yari intersequanta.}
\l{\cel{3}- DEFIS.}
\l{}
\l{\sou{perfekta}. Di qua la ecelo en sua genero esas absoluta.-DEFIS}
\l{}
\l{\sou{perfekto}. (\sou{gram.)}\cel{2}Tempo di verbo, qua indikas, ye maniero}
\l{\cel{3}absoluta, ke ago eventis, sen dicar precize ye qua epoko,}
\l{\cel{3}ye qua instanto. - DEFIS.}
\l{}
\l{\sou{perfida}. - Qua agas quale trahizanto, desloyale, ad ulu qua}
\l{\cel{3}fidis lu. - DEFIS.}
\l{}
\l{\sou{perforar}. (\sou{trans.)}\cel{2}Trairar facante truo, ma sen rotacar.}
\l{\cel{3}DEFIS.}
\l{}
\l{\sou{pergameno}. Felo de mutono, preparita por recevar skriburo,}
\l{\cel{3}imprimuro, pikturo, desegnuro. - DEFIRS.}
\l{}
\l{\sou{perianto.}\cel{2}(\sou{bot.)}\cel{1}Ensemblo ek la envelopili di floro, qua po-}
\l{\cel{3}vas konsistar ye kalico e korolo. - DEFIS.}
\l{}
\l{\sou{periferio}. I. (\sou{geom.)}\cel{3}Konturo di figuro kurvoforma. - II.}
\l{\cel{3}Surfaco extera di korpo. - DEFIRS.}
\l{}
\l{\sou{perifrazo}.\cel{2}(\sou{retoriko}). Frazo-formo uzata kom equivalanto di}
\l{\cel{3}la vorto apta, konvenanta. - DEFIRS. (\sou{perifrazar}).- DEFIRS.}
\l{}
\l{\sou{per}igeo. (astr.)\cel{2}Punto ek la orbito di planeto ube\cel{1}\sou{olca}\cel{1}esas}
\l{\cel{3}maxime proxima a Tero. - DEFIRS.}
}
